% Optional full-length-manuscript insert. This is intentionally not inserted
% into the ApJL draft: the zero-dof inversion and cross-sample He II gate need
% the methodological space and caveats retained below.

\section{A spectroscopic test of the physical visibility phase}
\label{sec:spectral_visibility}

The demographic model requires an observation layer that predicts more than
$n(z)$. We therefore represent the line-emitting nucleus as two physically
distinct photoionized zones: a dense compact skin, in which collisionally
excited forbidden lines are suppressed, and a lower-density channel or skin
that produces the optical metal lines. For line $\ell$ the emergent energy
yield per incident hydrogen-ionizing photon is
\begin{equation}
 y_\ell=f_{\rm d}\,y_{\ell,{\rm dense}}
 +(1-f_{\rm d})\,y_{\ell,{\rm diffuse}},
\end{equation}
where $f_{\rm d}$ is the fraction of incident $Q({\rm H})$ intercepted by the
dense phase. It is not a mass fraction, volume filling factor, or observed
broad-line fraction. Ratios are formed only after mixing the line yields.
The dense and diffuse yields are calculated in thermal balance with Cloudy
C25 at $Z=0.2Z_\odot$; positive yields are interpolated in $\log U$ within the
computed grid.

We confront this layer with three public constraints. First, the RUBIES
release contains 36 spectroscopic LRDs and 80 broad-Balmer sources
\citep{hviding2025}; 17 LRDs lie at $4.5\le z\le6.5$. Second,
\citet{perezgonzalez2026} report four stacked spectra from 249 LRDs, including
[O~III]~$\lambda5007$/H$\beta$, [O~III]/[O~II], H$\alpha$/H$\beta$, and
H$\beta$/H$\gamma$. Third, \citet{sok2026} find that most objects in their
high-S/N grating sample have He~II~$\lambda4686$/H$\beta\lesssim0.1$, while
low-density channels are required in at least some LRDs. Because these samples
have different selections, the He~II result is used as a population
compatibility test rather than as four subtype-specific measurements.

\paragraph{The ionizing spectrum must be soft above the He$^+$ edge.}
The original standard-AGN anchors predict He~II/H$\beta=0.244$ in the diffuse
zone and 0.729 in the dense zone, inconsistent with the typical observed
upper envelope when $f_{\rm d}$ is large. Thermal hardness-basis calculations
show that a 50-kK blackbody shape gives 0.0012 and 0.0186, respectively, while
still producing strong [O~III]. The temperature is not interpreted as a
literal source temperature; the inference is that the spectrum reaching the
gas must retain enough 35--54 eV photons to make O$^{++}$ while suppressing
photons above the 54.4-eV He$^+$ edge. A two-stage calculation in which a
standard AGN continuum crosses a generic ionized screen leaves He~II nearly
unchanged, so ordinary warm absorption is not sufficient in the tested
family.

\paragraph{The stacks define a dense-to-diffuse visibility sequence.}
For each stack, [O~III]/H$\beta$ and [O~III]/[O~II] determine $f_{\rm d}$ and
the diffuse $\log U$. With the He~II-compatible soft basis, the inferred
$f_{\rm d}$ values are $0.981^{+0.002}_{-0.003}$, $0.901^{+0.011}_{-0.017}$,
$0.738^{+0.069}_{-0.113}$, and $0.086^{+0.572}_{-0.086}$ for the xLRD,
plusLRD, minusLRD, and bLRD stacks, while $\log U$ remains near
$-2.20$ to $-2.08$. The bLRD fraction is weakly constrained, but the three
redder stacks show an ordered sequence. This supports the interpretation that
spectral diversity is primarily a changing dense intercepted-photon fraction
within a common stratified structure, rather than evidence that each subtype
requires a separate formation route.

\begin{figure}[t]
\centering
\includegraphics[width=\linewidth]{public_lrd_constraints/public_lrd_stack_two_zone_fits.pdf}
\caption{Two-zone inversion of the four LRD stacks. Left: measured ratios
(circles with uncertainties) and the corresponding fitted points (crosses)
within the Cloudy mixing sequences. Middle: the independently predicted
He~II/H$\beta$ ratios remain below the typical observed upper envelope for the
preferred soft basis. Right: observed Balmer ratios and the best one-parameter
smooth-screen predictions.}
\label{fig:lrd_stack_inversion}
\end{figure}

This inversion is not by itself a goodness-of-fit test: two ratios determine
two parameters. Its value is that the independent He~II constraint rejects
harder continua and that additional lines now provide over-identifying
falsifiers. In particular, He~I~$\lambda5876/7065/10830$, [Ne~III]/[O~II],
[O~I]/H$\alpha$, and [O~III]~$\lambda4363$ must be predicted jointly rather
than fitted one at a time.

\paragraph{Balmer transfer remains degenerate.}
The intrinsic two-zone Balmer decrements are close to recombination values,
whereas the stacks reach H$\alpha$/H$\beta=5.5$--10.1$. A one-parameter
power-law foreground screen gives $\chi^2=4.52$ for four nominal degrees of
freedom ($p=0.34$) and is therefore not rejected by the current
H$\gamma$ uncertainties. It requires effective $A_V\simeq2.1$--3.7 mag
toward the Balmer-emitting region. Optically thick line transfer and
differential broad/narrow attenuation remain alternatives, not demonstrated
necessities. A porous geometry can expose a blue continuum while obscuring
the line-emitting region, but this must be tested with a joint
continuum-plus-line likelihood.

\paragraph{An independent compact scale.}
Evaluating virial widths at the 20--40 pc nuclear-reservoir radius predicts a
median broad-Balmer FWHM of only $53\,{\rm km\,s^{-1}}$ for the redshift-matched
lifecycle population. A distinct compact line-emitting radius gives the
observed scale: median matching to the 17 RUBIES LRDs yields approximately
1220 au and a model median of $2246\,{\rm km\,s^{-1}}$, compared with the
observed $2245\,{\rm km\,s^{-1}}$. This is a descriptive calibration, not a
radius posterior, because the published Gaussian fits impose an upper width
bound of $2500\,{\rm km\,s^{-1}}$.

Together, these tests favor revising the visibility and ionizing-spectrum
calibration before invoking separate formation routes for the spectral
subtypes. They do not remove the independent $z\gtrsim8$ demographic tension,
which can still reflect an early formation/luminosity channel, revised
completeness, or both.

% Bibliography keys to add:
% hviding2025: Hviding et al. 2025, RUBIES LRD tables, Zenodo 16758549.
% perezgonzalez2026: Perez-Gonzalez et al. 2026, arXiv:2602.20247.
% sok2026: Sok et al. 2026, arXiv:2606.23778.
